% Part: methods
% Chapter: induction
% Section: introduction

\documentclass[../../../include/open-logic-section]{subfiles}

\begin{document}

\olfileid[hi]{mth}{ind}{int}

\olsection{परिचय}

आगमन एक महत्त्वपूर्ण प्रमाण तकनीक है जिसका विभिन्न रूपों में तर्कशास्त्र,
सैद्धांतिक कंप्यूटर विज्ञान और गणित के लगभग सभी क्षेत्रों में उपयोग होता
है। तर्कशास्त्र के अनेक परिणाम सिद्ध करने के लिए इसकी आवश्यकता पड़ती है।

आगमन की तुलना प्रायः निगमन से की जाती है और उसे विशिष्ट से सामान्य की ओर
अनुमान कहा जाता है। उदाहरण के लिए यदि हम अनेक हरे पन्ने देखते हैं और ऐसी
कोई वस्तु नहीं देखते जिसे हम पन्ना कहें पर जो हरी न हो, तो निष्कर्ष निकाल
सकते हैं कि सभी पन्ने हरे हैं। यह आगमनात्मक अनुमान है, क्योंकि यह अनेक
विशिष्ट स्थितियों (यह पन्ना हरा है, वह पन्ना हरा है, इत्यादि) से सामान्य
दावे (सभी पन्ने हरे हैं) की ओर जाता है। \emph{गणितीय} आगमन भी ऐसा अनुमान है
जो सामान्य दावा निष्कर्षित करता है, किंतु वह इस ``सरल आगमन'' से बहुत भिन्न
प्रकार का है।

बहुत मोटे तौर पर गणित का आगमनात्मक प्रमाण निष्कर्ष निकालता है कि किसी
विशेष प्रकार की सभी गणितीय वस्तुओं में कोई विशेष गुण है। सरलतम स्थिति में
आगमनात्मक प्रमाण जिन गणितीय वस्तुओं से संबंधित होता है वे प्राकृतिक संख्याएँ
हैं। उस स्थिति में आगमनात्मक प्रमाण यह स्थापित करता है कि सभी प्राकृतिक
संख्याओं में कोई गुण है, और इसके लिए वह दिखाता है कि
\begin{enumerate}
    \item $0$ में वह गुण है, और
    \item जब भी किसी संख्या $k$ में वह गुण हो, $k+1$ में भी वह गुण है।
\end{enumerate}
तब प्राकृतिक संख्याओं पर आगमन का उपयोग प्रायः उन गणितीय वस्तुओं के बारे में
सामान्य दावे सिद्ध करने में भी हो सकता है जिन्हें संख्याएँ दी जा सकती हैं।
उदाहरणतः प्रत्येक परिमित समुच्चय में अवयवों की कोई परिमित संख्या $n$ होती
है, और यदि आगमन से दिखा सकें कि प्रत्येक संख्या $n$ में यह गुण है कि
``$n$ आकार के सभी परिमित समुच्चय \dots हैं'', तो हमने सभी परिमित समुच्चयों
के बारे में कुछ दिखा दिया होगा।

आगमन का सामान्यीकरण उन गणितीय वस्तुओं तक भी किया जा सकता है जो
\emph{आगमनात्मक रूप से परिभाषित} हैं। उदाहरणतः प्रथम-क्रम तर्कशास्त्र जैसी
औपचारिक भाषा की अभिव्यक्तियाँ आगमनात्मक रूप से परिभाषित होती हैं।
\emph{संरचनात्मक आगमन} ऐसी सभी अभिव्यक्तियों के बारे में परिणाम सिद्ध करने
का ढंग है। विशेषतः संरचनात्मक आगमन तर्कशास्त्र में बहुत उपयोगी---और व्यापक
रूप से प्रयुक्त---है।

\end{document}
